% ------------------------------------------------------------------
% Sección 6.1: Síntesis de Hallazgos
% ------------------------------------------------------------------
\section{Síntesis de Hallazgos}
\label{sec:cap6_sintesis}

Esta tesis doctoral desarrolló, validó y operacionalizó un sistema híbrido integrado por sensores IoT de bajo costo, un modelo matemático mecanicista de balance de carbono (Dynamic Energy Budget) y un calibrador neuronal paramétrico, para la estimación en tiempo real de emisiones de CO$_2$ y la detección de eventos de anaerobiosis (CH$_4$) en la bioconversión de residuos agroindustriales con \emph{Hermetia illucens}. A continuación se sintetizan los aportes por capítulo y su articulación en un cuerpo de conocimiento coherente.

\textbf{Capítulo 1 -- Introducción.} Se situó la bioconversión con BSF como alternativa de economía circular para la gestión de residuos orgánicos en contextos tropicales, identificando el vacío crítico en la cuantificación dinámica de emisiones de gases de efecto invernadero. Se formuló la pregunta de investigación, las hipótesis (H1--H4) y los objetivos específicos que orientaron el desarrollo del sistema híbrido.

\textbf{Capítulo 2 -- Instrumentación y Metrología.} Se diseñó, implementó y caracterizó un módulo de adquisición basado en sensores NDIR de bajo costo (MH-Z19B para CO$_2$, MQ-4 para CH$_4$) integrado en un nodo IoT con procesamiento en el borde (ESP32). El protocolo de ventanas cerradas de 30 min, con intercalación de fases de aireación, permitió cuantificar tasas de emisión sin perturbar la oxigenación del reactor. La calibración frente a patrones certificados y la compensación cruzada por temperatura y humedad relativa redujeron el MAPE del sensor NDIR del 18{,}4\% (crudo) al 9{,}8\% (corregido), con correlación de Pearson $R = 0{,}997$ respecto al patrón y una incertidumbre expandida del 4{,}9\% para concentraciones de CO$_2$. Estos valores satisfacen holgadamente los umbrales de la hipótesis H1 ($R > 0{,}85$, MAPE $< 15$\%) y validan la viabilidad de sensores de bajo costo para monitoreo continuo en entornos agroindustriales tropicales.

\textbf{Capítulo 3 -- Modelo Matemático DEB.} Se desarrolló la generalización del modelo Dynamic Energy Budget de Eriksen (2022) desde la escala individual larval hasta la escala de reactor batch de 20~L, formulando un sistema de cinco ecuaciones diferenciales ordinarias que gobiernan la dinámica del sustrato orgánico, la reserva de asimilados, la biomasa estructural, los lípidos de reserva y la acumulación de CO$_2$. Se verificó la consistencia dimensional de todos los flujos (g~C/min), el cierre del balance de carbono en régimen aeróbico y el comportamiento cualitativo biológicamente plausible: no negatividad del dominio $\Omega$, monotonicidad decreciente del sustrato y acotación superior de la biomasa por el carbono total disponible. Se identificó el subconjunto de parámetros de primera sensibilidad respecto a la pendiente de emisión de CO$_2$: $\{\tau_h, \rho_{\mathrm{conv}}, T_{\mathrm{opt}}\}$.

\textbf{Capítulo 4 -- Arquitectura y Construcción del Sistema Híbrido.} Se formuló y validó la arquitectura de seis bloques funcionales (DEB Forward, sensor NDIR, extractor de pendientes, comparador, MLP calibrador y estimador de estado) que opera en lazo cerrado de calibración discreta. El MLP de dos capas ocultas (16+8 neuronas) actúa como corrector de escala sobre el subconjunto de primera sensibilidad, con salida acotada ($\pm 20$\%) para preservar la interpretabilidad biofísica. La estrategia de calibración por fases (semillas Eriksen, ajuste local L-BFGS-B, corrección neuronal en tiempo real) preserva la interpretabilidad del núcleo mecanicista mientras adapta el modelo a la variabilidad inter-réplica. Críticamente, en la Sección~\ref{sec:cap4_comparativa} se demostró mediante validación cruzada por réplica que una red neuronal directa (RNA-B) que predice biomasa desde variables operativas falla sistemáticamente (MAPE de validación $> 40$\%, predicciones no físicas), justificando formalmente la adopción del calibrador paramétrico.

\textbf{Capítulo 5 -- Validación y Diagnóstico Operativo.} El sistema híbrido validado predice la tasa de emisión de CO$_2$ con MAPE $= 8{,}3$\% y $R^2 = 0{,}78$, supera ampliamente el modelo lineal estático de referencia, detecta eventos de anaerobiosis con recall del 100\% (4 de 4 eventos confirmados) y cuantifica la discrepancia de carbono con trazabilidad metrológica ($\Delta_{\mathrm{Total}} = +5{,}1$~g~C, 3{,}6\% del total). La plataforma de soporte a la decisión, con arquitectura de tres capas (percepción, procesamiento, aplicación), traduce las salidas del modelo en un semáforo de riesgo ambiental (verde/ámbar/rojo) y un protocolo de respuesta ante anaerobiosis con tiempos cuantificados (confirmación 0--2~min, preparación 2--5~min, volteo 5--15~min, verificación 15--30~min). La generación automatizada de reportes MRV se alinea con estándares internacionales (IPCC Tier 3, Verra VCS v4.7), y las métricas de utilidad operativa validan la transición del prototipo (TRL 5) hacia una tecnología validada en entorno operativo (TRL 6).

\textbf{Aporte integrado.} El principal aporte de esta tesis es demostrar que la integración de sensores IoT de bajo costo, modelos mecanicistas bioenergéticos y calibradores neuronales paramétricos permite transformar la bioconversión con BSF de un proceso de ``caja negra'' en un sistema transparente, tecnológicamente viable y verificable climáticamente. La predicción matemática, la detección de anomalías y la cuantificación de carbono se materializan en decisiones de manejo, registros auditables y reportes comparables con estándares globales, fortaleciendo la autonomía técnica del sur global frente a los mercados de carbono y las políticas de economía circular.
